\documentclass[11pt,a4paper]{article}

\usepackage[margin=1in]{geometry}
\usepackage{times}
\usepackage{titlesec}
\usepackage{enumitem}
\usepackage{hyperref}
\usepackage{cite}
\usepackage{booktabs}
\usepackage{array}
\usepackage{parskip}

\hypersetup{
    colorlinks=true,
    linkcolor=black,
    citecolor=black,
    urlcolor=blue
}

\titleformat{\section}{\large\bfseries}{\thesection}{1em}{}
\titleformat{\subsection}{\normalsize\bfseries}{\thesubsection}{1em}{}
\titlespacing{\section}{0pt}{12pt}{6pt}
\titlespacing{\subsection}{0pt}{10pt}{4pt}

\setlength{\parindent}{0pt}
\setlength{\parskip}{6pt}

\title{\vspace{-1.5cm}\textbf{Early Sepsis Prediction from ICU Time Series:\\
Does Modelling Missingness Patterns Improve Prediction?}\\[4pt]
\large Project Synopsis}

% ============================================================
% EDIT THIS BLOCK: replace with your group's actual names/CMS IDs.
% Keep whichever line applies and delete the other two placeholder
% options, or leave all three if your template requires a fixed
% three-member format.
% ============================================================
\author{
\begin{tabular}{lll}
\textbf{Name} & \textbf{CMS ID} & \textbf{Role} \\
\hline
Student Name 1 & 000000 & \\
Student Name 2 & 000000 & \\
Student Name 3 & 000000 & \\
\end{tabular}
}

\date{\today}

\begin{document}

\maketitle
\thispagestyle{empty}
\vspace{-0.5cm}

\noindent\textbf{Course:} Machine Learning \hfill \textbf{Submission Date:} September 25, 2026

\vspace{0.3cm}
\hrule
\vspace{0.4cm}

% ============================================================
\section{Motivation and Problem Statement}

Sepsis is a leading cause of intensive care unit (ICU) mortality, and delays in
treatment are strongly associated with worse patient outcomes. Early sepsis,
however, presents with subtle, easily overlooked physiological changes that
closely resemble normal clinical variation, making early detection difficult
even for experienced clinicians. Automated early-warning systems built on ICU
time-series data offer a way to flag at-risk patients hours before clinical
recognition, potentially enabling earlier intervention.

A central challenge in this setting is that ICU data exhibits severe and
\emph{non-random} missingness: laboratory tests are ordered selectively based
on clinical suspicion rather than at fixed intervals, so the presence or
absence of a measurement is itself informative. Most existing approaches treat
missing data purely as an obstacle to be imputed away, discarding this
signal. This project investigates whether explicitly modelling missingness
patterns, rather than merely filling them in, improves early sepsis
prediction.

% ============================================================
\section{Related Work and Positioning}

Lipton et al.~\cite{lipton2016} demonstrated that treating missingness
indicators as explicit features, rather than relying on imputation alone,
improves predictive performance in clinical time-series classification. Che
et al.~\cite{che2018} extended this idea with the GRU-D architecture, formally
identifying \emph{informative missingness} as an underexploited signal in
multivariate clinical time series. Both works, however, rely on deep
recurrent architectures whose internal handling of missingness is difficult
to interpret in isolation.

Within the PhysioNet/CinC 2019 Sepsis Challenge specifically, the majority of
published approaches use gradient-boosted tree ensembles (e.g., XGBoost) with
imputation strategies such as mean-filling, linear interpolation, or
MissForest~\cite{cinc2019survey}. Where missingness indicators are used, they
are typically bundled together with many other engineered features without
isolating their individual contribution to performance.

This project differs by running a \textbf{controlled comparison} on the same
classifiers across three explicit data representations: (A) standard
imputation only, (B) imputation plus binary missingness masks, and (C)
imputation plus time-since-last-measurement features. Holding the model
architecture fixed while varying only the missingness representation isolates
the specific contribution of missingness modelling to prediction quality, a
comparison that is largely absent from prior work on this dataset.

% ============================================================
\section{Dataset}

We use the PhysioNet/Computing in Cardiology Challenge 2019
dataset~\cite{reyna2019}, comprising hourly ICU records for 40{,}336 patients
across two hospital systems (Hospital A: 20{,}336 patients; Hospital B:
20{,}000 patients), totaling 1{,}552{,}210 patient-hours. Each hourly record
contains 8 vital signs, 26 laboratory values, and 6 demographic variables,
with a binary \texttt{SepsisLabel} marking sepsis onset per Sepsis-3 criteria,
shifted six hours earlier than clinical recognition to reward early
prediction.

Exploratory analysis of the full dataset confirms the challenges motivating
this project. Class imbalance is severe at the row level: 7.27\% of patients
ever develop sepsis, but only 1.80\% of individual hourly rows are labeled
positive. Missingness is pronounced and highly uneven across features, as
summarized in Table~\ref{tab:missingness}.

\begin{table}[h]
\centering
\caption{Missingness in the five worst-affected laboratory features.}
\label{tab:missingness}
\begin{tabular}{lc}
\toprule
\textbf{Feature} & \textbf{\% Missing} \\
\midrule
Bilirubin\_direct & 99.81 \\
Fibrinogen        & 99.34 \\
TroponinI         & 99.05 \\
Bilirubin\_total   & 98.51 \\
Alkalinephos      & 98.39 \\
\bottomrule
\end{tabular}
\end{table}

This level of missingness is consistent with selective, clinically-driven
test ordering rather than random sensor dropout, making the dataset directly
suited to studying informative missingness. Its open access, established
provenance, and manageable scale make it feasible to work with within the
course timeframe.

% ============================================================
\section{Performance Metrics}

Given the severe row-level class imbalance (98.2\% negative), standard
accuracy is uninformative: a trivial classifier that always predicts
``no sepsis'' would score 98.2\%. We therefore report:

\begin{itemize}[itemsep=2pt]
    \item \textbf{PR-AUC} (precision--recall area under curve), which is more
    sensitive to minority-class performance than ROC-AUC under severe
    imbalance;
    \item \textbf{F1-score} at a selected operating threshold, balancing
    precision and recall;
    \item The Challenge's own \textbf{time-dependent utility score}, which
    rewards predictions made between 12 hours before and 3 hours after true
    sepsis onset, and penalizes both late detection and false alarms.
\end{itemize}

The utility score is essential to this project because it captures the
\emph{clinical} cost structure of an early-warning system in a way that
PR-AUC and F1 do not: a model can achieve strong discriminative performance
while still alerting too late to be clinically useful. Comparing
representations A, B, and C across all three metrics will indicate whether
missingness-aware modelling improves not only raw discrimination, but the
clinically-relevant early-detection property the utility score is designed to
capture.

% ============================================================
\begin{thebibliography}{9}

\bibitem{lipton2016}
Z.~C. Lipton, D.~C. Kale, and R.~Wetzel, ``Modeling Missing Data in Clinical
Time Series with RNNs,'' \emph{arXiv preprint arXiv:1606.04130}, 2016.

\bibitem{che2018}
Z.~Che, S.~Purushotham, K.~Cho, D.~Sontag, and Y.~Liu, ``Recurrent Neural
Networks for Multivariate Time Series with Missing Values,'' \emph{Scientific
Reports}, vol.~8, no.~6085, 2018.

\bibitem{reyna2019}
M.~A. Reyna, C.~S. Josef, R.~Jeter, S.~P. Shashikumar, M.~B. Westover,
S.~Nemati, A.~Sharma, and G.~D. Clifford, ``Early Prediction of Sepsis From
Clinical Data: The PhysioNet/Computing in Cardiology Challenge,''
\emph{Critical Care Medicine}, vol.~48, no.~2, pp.~210--217, 2019.

\bibitem{cinc2019survey}
Survey of PhysioNet/CinC 2019 Challenge entries (e.g., XGBoost-based
pipelines using mean imputation, linear interpolation, and MissForest for
missing-value handling), as reported in Challenge proceedings papers,
Computing in Cardiology, 2019.

\end{thebibliography}

\end{document}
